% Kapitel 6 — Budget: 10 Seiten — Schreiben: 19.10.–13.11.;
% 30.11.–11.12. auf den gefahrenen Stand bringen (Generalprobe, Pilot, Sitzungen).
% Bögen, Prüflisten und Leitfaden in den Anhang.

\chapter{Studiendesign}
\label{chap:studiendesign}

\section{Fragen der Studie und ihre Herkunft aus den Vormessungen}% 1,5 S.
\label{sec:studienfragen}

\section{Bedingungen C und D}% 1 S.
\label{sec:bedingungen}

\section{Aufgaben auf teamable}% 2 S.
% Inhalt: Referenzkontext, Spielraum je Aufgabe (13.10.).
\label{sec:aufgaben}

\section{Messung}% 2,5 S.
% Inhalt: Vorfragebogen, Fragen vorher und nachher, Erklären, UI-Export,
% Prüfliste für das Ergebnis des Agenten, TLX, SUS, Interview.
\label{sec:messung}

\section{Teilnehmende, Ablauf, Einwilligung, Pilot und Generalprobe}% 1,5 S.
\label{sec:teilnehmende}

\section{Auswertung}% 1,5 S.
% Inhalt: beschreibend und explorativ angelegt, acht Differenzen, Vorzeichentest
% ergänzend, zwei Schichten (geliefert, nicht geliefert), Kontrollitems gegen den
% Vortest als Suchauftrag, Entwickler getrennt.
\label{sec:auswertung}
